% part: intuitionistic-logic
% chapter: introduction
% section: constructive-reasoning

\documentclass[../../../include/open-logic-chapter]{subfiles}

\begin{document}

\olfileid{int}{int}{cr}

\section{నిర్మాణాత్మక హేతుచింతన}

సాంప్రదాయిక తర్కాన్ని మోడల్ కారకాలతో లేదా
ద్వితీయ-స్థాయి పరిమాణీకరణికాలతో విస్తరించడానికి
భిన్నంగా, అంతఃప్రజ్ఞావాద తర్కం సాంప్రదాయిక
తర్కాన్ని పరిమితం చేస్తుంది కాబట్టి ``అసాంప్రదాయికం''.
సాంప్రదాయిక తర్కం అనేక విధాలుగా
\emph{నిర్మాణాత్మకం కానిది}. ఒక రకమైన నిర్మాణాత్మక
గణితానికి లక్షణమైన మరింత ``నిర్మాణాత్మక''
హేతుచింతనను పట్టుకోవడానికి అంతఃప్రజ్ఞావాద తర్కాన్ని
ఉద్దేశించారు. దాని వెనుక ప్రేరణల్లో కొన్నింటిని
కింది ఉదాహరణలు వివరించవచ్చు.

ఎవరైనా కింది ధర్మం గల సహజ సంఖ్య~$n$ను
నిర్ణయించానని ప్రకటించారనుకోండి: $n$ సరి సంఖ్య
అయితే రీమాన్ పరికల్పన సత్యం; $n$ బేసి సంఖ్య
అయితే రీమాన్ పరికల్పన అసత్యం. అద్భుతమైన వార్త!{}
రీమాన్ పరికల్పన సత్యమా కాదా అనేది గణితంలోని
పెద్ద అపరిష్కృత ప్రశ్నల్లో ఒకటి; ఇక్కడ వారు
సమస్యను ఒక గణనకు, అంటే నిర్దిష్ట సంఖ్య సరిదా
కాదా అని నిర్ణయించడానికి తగ్గించినట్లు కనిపిస్తోంది.

$n$ యొక్క ఆ అద్భుత విలువ ఏమిటి? వారు దాన్ని
ఇలా వివరిస్తారు: రీమాన్ పరికల్పన సత్యమైతే
$2$కు, కాకపోతే $3$కు సమానమైన సహజ సంఖ్యే~$n$.

కోపంతో మీరు చెల్లించిన డబ్బు తిరిగి ఇవ్వాలని
అడుగుతారు. సాంప్రదాయిక దృక్కోణం నుంచి
పైనున్న వివరణ నిజంగానే $n$కు ఒకే విలువను
నిర్ణయిస్తుంది; కానీ మీకు నిజంగా కావలసింది
\emph{స్పష్టంగా} ఇచ్చిన $n$ విలువ.

మరొక, బహుశా తక్కువ కల్పితమైన ఉదాహరణగా
కింది ప్రశ్నను పరిశీలించండి. ఒక అకరణీయ సంఖ్యను
కరణీయ ఘాతానికి హెచ్చించి కరణీయ ఫలితాన్ని
పొందడం సాధ్యమని మనకు తెలుసు. ఉదాహరణకు,
$\sqrt{2}^2 = 2$. ఒక అకరణీయ సంఖ్యను
\emph{అకరణీయ} ఘాతానికి హెచ్చించి కరణీయ
ఫలితాన్ని పొందడం సాధ్యమా అనేది అంత స్పష్టం
కాదు. కింది సిద్ధాంతం దీనికి అవునని సమాధానం
ఇస్తుంది:

\begin{thm}
$a^b$ కరణీయమయ్యే అకరణీయ సంఖ్యలు $a$, $b$ ఉన్నాయి.
\end{thm}

\begin{proof}
$\sqrt{2}^{\sqrt{2}}$ను పరిశీలించండి. ఇది
కరణీయమైతే మన పని పూర్తయింది: $a = b = \sqrt{2}$గా
తీసుకోవచ్చు. కాకపోతే అది అకరణీయం. అప్పుడు
మనకు
\[
(\sqrt{2}^{\sqrt{2}})^{\sqrt{2}} = \sqrt{2}^{\sqrt{2} \cdot
  \sqrt{2}} = \sqrt{2}^2 = 2,
\]
లభిస్తుంది; ఇది కరణీయమే. కాబట్టి ఈ సందర్భంలో
$a$ను $\sqrt{2}^{\sqrt{2}}$గా, $b$ను~$\sqrt 2$గా
తీసుకోండి.
\end{proof}

ఇది చెల్లుబాటైన నిరూపణేనా? అవునని చాలామంది
గణితశాస్త్రవేత్తలు భావిస్తారు. అయినా ఇక్కడ మళ్లీ
కొంత అసంతృప్తి కలిగించే విషయం ఉంది: ఒక
నిర్దిష్ట ధర్మం గల వాస్తవ సంఖ్యల జత ఉందని
నిరూపించాం; కానీ అది \emph{ఏ} జతో చెప్పలేకపోయాం.
అదే ఫలితాన్ని, $a$, $b$ జతను నిరూపణలోనే
\emph{ఇచ్చే} విధంగా నిరూపించవచ్చు:
$a = \sqrt{3}$, $b = \log_3 4$గా తీసుకోండి. అప్పుడు
\[
a^b = \sqrt{3}^{\log_3 4} = 3^{1/2 \cdot \log_3 4} = (3^{\log_3
  4})^{1/2} = 4^{1/2}= 2,
\]
ఎందుకంటే $3^{\log_3 x} = x$.

మొదటి నిరూపణలోని అడుగుల వంటివాటిని అనుమతించని
హేతుచింతనను పట్టుకోవడానికి అంతఃప్రజ్ఞావాద
తర్కాన్ని రూపొందించారు. $!A(x)$ను సంతృప్తిపరచే
$x$ ఉందని నిరూపించడం అంటే, రెండవ నిరూపణలోలాగా
ఒక నిర్దిష్ట~$x$ను, అది $!A$ను సంతృప్తిపరుస్తుందనే
నిరూపణను ఇవ్వాలి. $!A$ లేదా $!B$ నిలుస్తుందని
నిరూపించాలంటే, వాటిలో ఒకదాన్ని నిరూపించగలగాలి.

ఔపచారికంగా చెప్పాలంటే, సాంప్రదాయిక తర్కపు
\tetoken{ఒక వ్యుత్పత్తి}{!!a{derivation}} వ్యవస్థను
ఒక నిర్దిష్ట విధంగా పరిమితం చేస్తే వచ్చేదే
అంతఃప్రజ్ఞావాద తర్కం. గణిత దృక్కోణం నుంచి
ఇవి కేవలం ఔపచారిక నిగమన వ్యవస్థలు; కానీ
ఇప్పటికే గమనించినట్లు, ఒక రకమైన గణిత
హేతుచింతనను పట్టుకోవడానికి వీటిని ఉద్దేశించారు.
గణితం గురించి ఒక నిర్దిష్ట తాత్త్విక దృక్కోణం
(బ్రౌవర్ అంతఃప్రజ్ఞావాదం వంటిది) సమర్థించే
హేతుచింతన ఇదని తీసుకోవచ్చు; మరింత
``మూర్తమైన'', సంతృప్తికరమైన గణిత హేతుచింతన
(బిషప్ నిర్మాణవాద పద్ధతిలోనిది) అని తీసుకోవచ్చు;
ఔపచారిక వివరణ అనౌపచారిక ప్రేరణను నిజంగా
పట్టుకుంటుందా అని వాదించవచ్చు. కానీ మన
తాత్త్విక వైఖరులు ఏవైనా, అంతఃప్రజ్ఞావాద తర్కాన్ని
ఔపచారికంగా సమర్పించిన తర్కంగా అధ్యయనం చేయవచ్చు;
ఏ కారణాల వల్లనైనా, అనేకమంది గణిత తర్కశాస్త్రవేత్తలు
అలా చేయడాన్ని ఆసక్తికరంగా భావిస్తారు.

\end{document}
